\section{Domain and assignment dependencies in the Considerations comparison}
\label{sec:p08-n-b-009-02}

The original 1963 \emph{Semantical Considerations on Modal Logic} is
distinct from \emph{Semantical Analysis of Modal Logic I}. The following
supporting treatment uses the separately identified 1971 reprint of
\emph{Considerations}, printed pp.~63--72, and proves its domain and
closure comparisons explicitly \cite{Kripke1971ConsiderationsReprint}.
It is not a replacement of the exact 1963 source presentation: the
reprint adds A0 to the opening propositional list, and the original has
a footnote omitted in the reprint. The source's omitted tableaux and
reported completeness arguments are not represented as proofs supplied
by this semantic comparison.

\subsection{The reprint's vocabulary and two-valued convention}

\begin{hbdefinition}{def:p08-kc-model}{Pointed source structures, predicate extensions and assignments}
A source model structure is \((G,K,R)\), where K is a nonempty world
set, \(G\in K\) is the distinguished world and R is reflexive.
The source M is the reflexive normal system, the modern T, not Chellas's
monotonic M or Pacuit's EM. Adding transitivity gives its S4 structures;
adding symmetry gives the Brouwersche structures; equivalence relations
give its S5 structures. A quantificational model structure additionally
specifies \(\psi(H)\), a possibly empty set of individuals existing at
each \(H\in K\). There is no automatic inclusion condition along R.
Put \(U=\bigcup_{H\in K}\psi(H)\).

The selected language has individual variables, n-place predicate
letters for every finite n including zero, Boolean conjunction and
negation, \(\Box\), and the source universal binder \((x)\).
Implication and the other Boolean operations are abbreviations;
\(\Diamond A\defeq\neg\Box\neg A\).
There are no primitive individual constants or function symbols in
this selected quantificational syntax.
For n>0, \(\varphi(P^n,H)\subseteq U^n\); for n=0 its value
is T or F. Thus predicates may be true of tuples not all of whose
members exist at the evaluation world. This total two-valued policy
is one chosen convention, not a claim that competing truth-gap
conventions are incoherent.

For a formula A use an assignment \(s:\operatorname{FV}(A)\to U\).
Writing \(\varphi_s(A,H)\) makes the source's implicit parameter visible.
An atomic formula is true exactly when its assigned tuple belongs
to the indicated predicate extension. Boolean truth is classical,
modal movement keeps the same values of free variables, and
\[
 \varphi_s((x)A,H)=\mathrm T
 \quad\Longleftrightarrow\quad
 \varphi_{s[x\mapsto a]}(A,H)=\mathrm T
       \text{ for every }a\in\psi(H).
\]
Assignments are restricted to the subformula's free variables as needed.
The empty assignment evaluates a closed formula even if U is empty;
an open formula with free variables has no U-assignment in that case.
Quantifier truth is vacuous when its local domain is empty. This is
not the consumer's Fitting convention in which primitive \(\forall\)
ranges over U and \(\forall^E\) is a separate abbreviation
\cite[pp.~64--67]{Kripke1971ConsiderationsReprint}.
\end{hbdefinition}

\begin{hblemma}{thm:p08-kc-evaluation}{proved=annex:p8.kc.evaluation}{Total evaluation and capture-free variable replacement}
For each supplied assignment, every formula has a unique truth value
at every world, including worlds where an assigned individual is absent.
Evaluation depends only on its free-variable values.
If y is free for x, replacing free x by y agrees with evaluating the
original formula after assigning x the value assigned to y.
No step assumes that this value belongs to the current local domain.
\end{hblemma}

\subsection{Closed assertions and the necessitation obstruction}

\begin{hbdefinition}{def:p08-kc-closures}{Generality interpretation and arbitrary-order closures}
The reprint treats an assertion of A with free variables as an assertion
of its ordinary universal closure. It does not treat this abbreviation
as a theorem directly concerning every arbitrary U-assignment.
A source \emph{closure} prefixes universal quantifiers and necessity
operators, in any order, until the result has no free variables. The
ordering matters when local domains vary.

The quantified M axiom inventory consists of all such closures of:
truth-functional tautologies; \(\Box A\to A\);
\(\Box(A\to B)\to(\Box A\to\Box B)\);
\(A\to(x)A\) when x is not free in A;
\((x)(A\to B)\to((x)A\to(x)B)\); and
\((y)((x)A(x)\to A(y))\), with capture-free replacement.
Detachment on closed formulas is the sole primitive inference;
necessitation is derived. The source quantified extensions add all
closures of \(\Box A\to\Box\Box A\) for S4,
\(A\to\Box\Diamond A\) for the Brouwersche system, or
\(\Diamond A\to\Box\Diamond A\) for S5, respectively.
This description is theorem-level: no local-premise necessitation is
licensed, and no additional premise consequence relation is silently
specified \cite[pp.~68--70]{Kripke1971ConsiderationsReprint}.
\end{hbdefinition}

\begin{hbtheorem}{thm:p08-kc-closed-soundness}{proved=annex:p8.kc.closed-soundness}{Closed-axiom soundness and derived necessitation}
Every displayed closure schema is valid in the stated two-valued
reflexive-domain semantics. Closed detachment preserves validity and
necessitation is admissible on closed theorems by induction over their
proofs. Necessitating the \emph{assertion} of a free-variable formula
therefore boxes its already closed reading. It does not justify moving
its ordinary universal closure outside the newly added box.
\end{hbtheorem}

\begin{hbtheorem}{thm:p08-kc-necessitation-boundary}{proved=annex:p8.kc.necessitation-boundary}{The failed Barcan derivation's precise step}
The classical instantiation assertion
\((x)A(x)\to A(y)\) has closed reading
\((y)((x)A(x)\to A(y))\). Necessitation of that closed
theorem produces its boxed universal closure, not
\((y)\Box((x)A(x)\to A(y))\).
The latter transformation requires a converse-Barcan principle and
fails on the complete two-world S5 countermodel below. Subsequent
distribution and generalization steps do not repair this change of
assertion policy \cite[pp.~68--69]{Kripke1971ConsiderationsReprint}.
\end{hbtheorem}

\subsection{Barcan directions, exact countermodels and existence}

\begin{hbdefinition}{def:p08-kc-barcan}{The two named directions under the source binder}
Write
\[
 \mathrm{BF}:(x)\Box A(x)\to\Box(x)A(x),\qquad
 \mathrm{CBF}:\Box(x)A(x)\to(x)\Box A(x).
\]
These names stay in this order. Domains decrease along an edge
\(HRH'\) when \(\psi(H')\subseteq\psi(H)\); they increase
when the reverse inclusion holds. The assigned individuals remain
the same through either modal move, and predicates have the already
specified extensions on U, including absent-individual arguments.
Constant domains here mean constant on all worlds; both inclusions
on each edge would suffice for the two direction arguments.
\end{hbdefinition}

\begin{hbtheorem}{thm:p08-kc-barcan-domains}{proved=annex:p8.kc.barcan-domains}{Full direction proofs and converses with arbitrary predicates}
On any such frame, BF is valid under all predicate interpretations
exactly when domains decrease along every edge, and CBF exactly when
they increase along every edge. Thus both hold for constant domains;
increasing domains guarantee CBF but BF may fail, decreasing domains
guarantee BF but CBF may fail, and unconstrained domains permit both
failures. Empty local domains do not invalidate the positive direction
proofs. The converses use one unary predicate with arbitrary U-extension,
not a hidden restriction to currently existing arguments.
\end{hbtheorem}

\begin{hbtheorem}{thm:p08-kc-s5-countermodels}{proved=annex:p8.kc.s5-countermodels}{The two complete S5 witnesses in the reprint}
For both examples set \(K=\{G,H\}\), \(G\ne H\), and
\(R=K^2\), with distinct individuals a,b.
For BF take \(\psi(G)=\{a\}\), \(\psi(H)=\{a,b\}\),
and \(\varphi(P,G)=\varphi(P,H)=\{a\}\).
Then \((x)\Box P(x)\) is true at G but \(\Box(x)P(x)\)
false. For CBF take \(\psi(G)=\{a,b\}\), \(\psi(H)=\{a\}\),
\(\varphi(P,G)=\{a,b\}\), \(\varphi(P,H)=\{a\}\).
Then \(\Box(x)P(x)\) is true at G but \((x)\Box P(x)\)
false. These are not globally increasing/decreasing S5 models:
the reverse edge enforces the opposite inclusion. They demonstrate
failure of both schemas in quantified S5 with unrestricted domains
\cite[pp.~67--68]{Kripke1971ConsiderationsReprint}.
\end{hbtheorem}

\begin{hbdefinition}{def:p08-kc-existence}{Existence and the optional identity extension}
The source existence predicate has
\(\varphi(E,H)=\psi(H)\), as a unary extension on U.
The reprint introduces its axioms by closures of
\(((x)A(x)\land E(y))\to A(y)\) and \((x)E(x)\).
Its simple identity extension makes \(x=y\) true exactly when
the two U-assigned values are identical, regardless of current existence.
Under that equality, \(E(x)\) is definable as \((\exists y)(x=y)\).
The suggestion of broader identity semantics is a cited boundary,
not a completed alternative construction in this paper
\cite[p.~70]{Kripke1971ConsiderationsReprint}.
\end{hbdefinition}

\begin{hbtheorem}{thm:p08-kc-existence}{proved=annex:p8.kc.existence}{Identity, existence and unrestricted predicate substitution}
The stated existence axioms are sound and the quantified identity
definition gives exactly current-domain membership. \(\Box(x)E(x)\)
is valid, while \((x)\Box E(x)\) may fail. In contrast the
tautological predicate \(A(x)\lor\neg A(x)\) is true of every
U-assignment at every world, including absent individuals.

If one instead restricts every atomic predicate extension to
\(\psi(H)^n\), unrestricted substitution of arbitrary formulas for
predicates need not preserve the resulting valid closed formulas.
That is a distinct semantic/rule combination. The original total
two-valued branch does not impose this restriction merely because
it quantifies locally over existing objects.
\end{hbtheorem}

\subsection{Explicit later-model bridge and designation extension}

\begin{hbtheorem}{thm:p08-kc-fitting-bridge}{proved=annex:p8.kc.fitting-bridge}{The translated consumer comparison, without renaming the source quantifier}
When local domains are nonempty, use the same world set, R, U and
predicate extensions in the Fitting consumer of
Definition~\ref{def:p08-qml-model}. Replace each source \((x)A\)
by \((\forall^E x)t(A)\), keep modalities and Boolean operations,
and give the Fitting existence predicate extension \(\psi(H)\).
The translation preserves truth in the common variable/predicate/modal
language under every supplied assignment and at every world. Including
the optional identity language additionally requires the same U-value
identity clause in both models. The source binder is compared with Fitting's
\emph{relativized} binder, not its primitive possibilist \(\forall\).
The source's empty-domain cases have no image inside that particular
consumer's declared nonempty-domain class.

If instead every primitive quantifier ranges over fixed U, both
unrelativized Barcan directions are valid by commuting two universal
quantifications. Varying \(\psi\) then affects E but not those
unrelativized direction proofs. None of these translations turns a
philosophical actualist or possibilist position into a proved theorem.
\end{hbtheorem}

\begin{hbtheorem}{thm:p08-kc-designation}{proved=annex:p8.kc.designation}{Fixed variable values are not a complete theory of names}
In an explicitly added later-model constant extension, a rigid name
c denotes the same \(d\in U\) at all worlds, independently of whether
\(d\in\psi(H)\). Such a name may satisfy \(E(c)\) but fail
\(\Box E(c)\), while \(\Box(c=c)\) remains true.
If names vary with worlds, the source free-variable instantiation
argument cannot be reused for a name without a separate substitution
check. The annex supplies a full constant-domain countermodel to the
named instance of the fixed-variable necessary-identity principle;
it is not a counterexample to the Barcan domain theorem. This extension is not
misattributed to the variable-only syntax described above.
\end{hbtheorem}

\paragraph{The proof boundary.}
The reprint reports propositional completeness and extension of it to
the quantified systems, but explicitly omits tableaux and largely suppresses
proofs. Its closing provability interpretations are marked brief and sketchy;
they are not the later arithmetical provability semantics or an unrestricted
quantified completeness derivation. The complete original Acta manifestation
has its own version-sensitive allocation. The independently proved comparisons
here give no completeness credit to a missing original passage, and do not
discharge the separate Barcan, Fitting--Mendelsohn or Garson presentations.
